\documentclass[11pt]{article}
% The preamble every project's paper shares: % The preamble every project's paper shares: % The preamble every project's paper shares: \input{../../preamble} from
% projects/<project>/paper/. Boards come from chessboard and xskak, plots from
% pgfplots; every number a paper prints comes from its generated.tex.
\usepackage[LSBC4,T1]{fontenc}
\usepackage[utf8]{inputenc}
\usepackage{lmodern}
\usepackage[margin=1in]{geometry}
\usepackage{microtype}
\usepackage{amsmath}
\usepackage{booktabs}
\usepackage{tabularx}
\usepackage{array}
\usepackage{xcolor}
\usepackage[obeyspaces,spaces]{url}
\usepackage{xskak}
\usepackage{pgfplots}
\pgfplotsset{compat=1.18}
\usepackage{caption}
\usepackage{subcaption}
\usepackage{enumitem}
\usepackage[section]{placeins}
\usepackage[hidelinks]{hyperref}

\setchessboard{boardfontencoding=LSBC4, boardfontfamily=skaknew, blackfieldcolor=black!22,
  whitefieldcolor=white, setfontcolors, showmover=false, label=true, labelfontsize=6pt,
  boardfontsize=14pt}

% A move or line in UCI, and a position in FEN, both breakable at a slash.
\newcommand{\uci}[1]{\texttt{#1}}
\newcommand{\fen}[1]{{\footnotesize\path{#1}}}
\newcolumntype{L}{>{\raggedright\arraybackslash}X}

\setlength{\parskip}{0.4em}
\setlength{\parindent}{0pt}
 from
% projects/<project>/paper/. Boards come from chessboard and xskak, plots from
% pgfplots; every number a paper prints comes from its generated.tex.
\usepackage[LSBC4,T1]{fontenc}
\usepackage[utf8]{inputenc}
\usepackage{lmodern}
\usepackage[margin=1in]{geometry}
\usepackage{microtype}
\usepackage{amsmath}
\usepackage{booktabs}
\usepackage{tabularx}
\usepackage{array}
\usepackage{xcolor}
\usepackage[obeyspaces,spaces]{url}
\usepackage{xskak}
\usepackage{pgfplots}
\pgfplotsset{compat=1.18}
\usepackage{caption}
\usepackage{subcaption}
\usepackage{enumitem}
\usepackage[section]{placeins}
\usepackage[hidelinks]{hyperref}

\setchessboard{boardfontencoding=LSBC4, boardfontfamily=skaknew, blackfieldcolor=black!22,
  whitefieldcolor=white, setfontcolors, showmover=false, label=true, labelfontsize=6pt,
  boardfontsize=14pt}

% A move or line in UCI, and a position in FEN, both breakable at a slash.
\newcommand{\uci}[1]{\texttt{#1}}
\newcommand{\fen}[1]{{\footnotesize\path{#1}}}
\newcolumntype{L}{>{\raggedright\arraybackslash}X}

\setlength{\parskip}{0.4em}
\setlength{\parindent}{0pt}
 from
% projects/<project>/paper/. Boards come from chessboard and xskak, plots from
% pgfplots; every number a paper prints comes from its generated.tex.
\usepackage[LSBC4,T1]{fontenc}
\usepackage[utf8]{inputenc}
\usepackage{lmodern}
\usepackage[margin=1in]{geometry}
\usepackage{microtype}
\usepackage{amsmath}
\usepackage{booktabs}
\usepackage{tabularx}
\usepackage{array}
\usepackage{xcolor}
\usepackage[obeyspaces,spaces]{url}
\usepackage{xskak}
\usepackage{pgfplots}
\pgfplotsset{compat=1.18}
\usepackage{caption}
\usepackage{subcaption}
\usepackage{enumitem}
\usepackage[section]{placeins}
\usepackage[hidelinks]{hyperref}

\setchessboard{boardfontencoding=LSBC4, boardfontfamily=skaknew, blackfieldcolor=black!22,
  whitefieldcolor=white, setfontcolors, showmover=false, label=true, labelfontsize=6pt,
  boardfontsize=14pt}

% A move or line in UCI, and a position in FEN, both breakable at a slash.
\newcommand{\uci}[1]{\texttt{#1}}
\newcommand{\fen}[1]{{\footnotesize\path{#1}}}
\newcolumntype{L}{>{\raggedright\arraybackslash}X}

\setlength{\parskip}{0.4em}
\setlength{\parindent}{0pt}

% Written by projects/accuracy-vs-chesscom/paper/generate.mjs; do not edit by hand.

\newcommand{\Games}{60}
\newcommand{\GamesTrain}{30}
\newcommand{\GamesTest}{30}
\newcommand{\Eligible}{192}
\newcommand{\DevPool}{132}
\newcommand{\RejectedSplits}{1}
\newcommand{\PlayerRatingMin}{948}
\newcommand{\PlayerRatingMax}{1167}
\newcommand{\ShortMax}{40}
\newcommand{\MediumMin}{41}
\newcommand{\MediumMax}{80}
\newcommand{\WithinPoints}{5}
\newcommand{\RoutineScan}{1,000,000}
\newcommand{\PhoneScan}{100,000}
\newcommand{\VerifyNodes}{750,000}
\newcommand{\VerifyLines}{3}
\newcommand{\PhoneShare}{25}
\newcommand{\BaseTestMae}{7.5}
\newcommand{\BaseTestMaeCI}{7.5 [5.3, 10.0]}
\newcommand{\BaseTrainMae}{8.3}
\newcommand{\BaseTestWithin}{50}
\newcommand{\BaseTestBias}{$-$0.5}
\newcommand{\BaseTestR}{0.49}
\newcommand{\BaseTestRCI}{0.49 [0.29, 0.73]}
\newcommand{\OneTestMae}{4.9}
\newcommand{\OneTestMaeCI}{4.9 [3.7, 6.2]}
\newcommand{\OneTrainMae}{3.6}
\newcommand{\OneTestWithin}{63}
\newcommand{\OneTestBias}{0.5}
\newcommand{\OneTestR}{0.76}
\newcommand{\OneTestRCI}{0.76 [0.61, 0.88]}
\newcommand{\OneVsBase}{$-$2.6 [$-$4.8, $-$0.9]}
\newcommand{\TwoTestMae}{5.0}
\newcommand{\TwoTestMaeCI}{5.0 [3.8, 6.3]}
\newcommand{\TwoTrainMae}{3.6}
\newcommand{\TwoTestWithin}{63}
\newcommand{\TwoTestBias}{0.6}
\newcommand{\TwoTestR}{0.76}
\newcommand{\TwoTestRCI}{0.76 [0.61, 0.88]}
\newcommand{\TwoVsBase}{$-$2.5 [$-$4.7, $-$0.7]}
\newcommand{\PhoneBaseTestMae}{7.5}
\newcommand{\PhoneBaseTestMaeCI}{7.5 [5.5, 9.8]}
\newcommand{\PhoneBaseTrainMae}{8.1}
\newcommand{\PhoneBaseTestWithin}{48}
\newcommand{\PhoneBaseTestBias}{0.2}
\newcommand{\PhoneBaseTestR}{0.50}
\newcommand{\PhoneBaseTestRCI}{0.50 [0.30, 0.74]}
\newcommand{\PhoneOneTestMae}{5.4}
\newcommand{\PhoneOneTestMaeCI}{5.4 [4.0, 6.8]}
\newcommand{\PhoneOneTrainMae}{4.6}
\newcommand{\PhoneOneTestWithin}{57}
\newcommand{\PhoneOneTestBias}{0.0}
\newcommand{\PhoneOneTestR}{0.72}
\newcommand{\PhoneOneTestRCI}{0.72 [0.54, 0.86]}
\newcommand{\PhoneOneVsBase}{$-$2.1 [$-$4.1, $-$0.5]}
\newcommand{\PhoneTwoTestMae}{5.4}
\newcommand{\PhoneTwoTestMaeCI}{5.4 [4.0, 6.8]}
\newcommand{\PhoneTwoTrainMae}{4.6}
\newcommand{\PhoneTwoTestWithin}{57}
\newcommand{\PhoneTwoTestBias}{0.0}
\newcommand{\PhoneTwoTestR}{0.72}
\newcommand{\PhoneTwoTestRCI}{0.72 [0.54, 0.86]}
\newcommand{\PhoneTwoVsBase}{$-$2.1 [$-$4.1, $-$0.5]}
\newcommand{\AggregateTable}{%
\begin{tabular}{lllllll}
\toprule
instrument & test MAE & train MAE & within 5 & bias & $r$ & vs baseline \\
\midrule
\multicolumn{7}{l}{\emph{routine evidence (the headline)}} \\
baseline (Lichess) & 7.5 [5.3, 10.0] & 8.3 & 50\% & $-$0.5 [$-$4.1, 2.8] & 0.49 [0.29, 0.73] & -- \\
round 1 & 4.9 [3.7, 6.2] & 3.6 & 63\% & $+$0.5 [$-$1.6, 2.6] & 0.76 [0.61, 0.88] & $-$2.6 [$-$4.8, $-$0.9] \\
the app (v2) & 5.0 [3.8, 6.3] & 3.6 & 63\% & $+$0.6 [$-$1.5, 2.7] & 0.76 [0.61, 0.88] & $-$2.5 [$-$4.7, $-$0.7] \\
\midrule
\multicolumn{7}{l}{\emph{phone evidence (the guardrail)}} \\
baseline (Lichess) & 7.5 [5.5, 9.8] & 8.1 & 48\% & $+$0.2 [$-$3.3, 3.3] & 0.50 [0.30, 0.74] & -- \\
round 1 & 5.4 [4.0, 6.8] & 4.6 & 57\% & $+$0.0 [$-$2.3, 2.4] & 0.72 [0.54, 0.86] & $-$2.1 [$-$4.1, $-$0.5] \\
the app (v2) & 5.4 [4.0, 6.8] & 4.6 & 57\% & $+$0.0 [$-$2.3, 2.4] & 0.72 [0.54, 0.86] & $-$2.1 [$-$4.1, $-$0.5] \\
\bottomrule
\end{tabular}
}
\newcommand{\ReproAggregates}{match the published ones to the character}
\newcommand{\NotRun}{%
\begin{itemize}[nosep]
\item book and forced moves counted: inside the noise or worse on the dev pool and train.
\item flatter cp-to-win curve for accuracy: inside the noise or worse on the dev pool and train.
\item floors and caps on per-move scores: 6.4 on dev, worse than round 1.
\item book moves as best on top of round 1: cut the length tilt, MAE inside noise.
\item per-category move scores (bands, classes, blunder cap): dev 4.46 / 4.50 / 4.30 against round 1's 4.29 in-sample; cross-validated 4.6.
\end{itemize}
}
\newcommand{\StopRule}{no candidate beat round 1 by more than 0.5 on the dev pool while improving train}
\newcommand{\ConstantGuess}{73.8}
\newcommand{\ConstantTestMae}{7.9}
\newcommand{\ConstantTrainMae}{7.2}
\newcommand{\TiltTable}{%
\begin{tabular}{lrrrr}
\toprule
length & sides & baseline & round 1 & the app (v2) \\
\midrule
short (40 plies or fewer) & 24 & $-$7.7 & $-$4.1 & $-$4.1 \\
medium (41 to 80) & 66 & $-$2.4 & $+$0.7 & $+$0.7 \\
long (over 80) & 30 & $+$5.1 & $+$3.5 & $+$3.9 \\
\bottomrule
\end{tabular}
}
\newcommand{\TiltBaseShort}{7.7}
\newcommand{\TiltTwoShort}{4.1}
\newcommand{\TiltBaseMedium}{2.4}
\newcommand{\TiltTwoMedium}{0.7}
\newcommand{\TiltBaseLong}{5.1}
\newcommand{\TiltTwoLong}{3.9}
\newcommand{\ScatterBase}{(84.52,87.61) (69.23,75.17) (70.61,66.84) (65.90,45.97) (82.44,80.32) (68.70,67.69) (77.37,77.42) (69.03,67.91) (71.07,67.51) (86.38,84.86) (75.99,91.89) (63.99,80.60) (62.50,66.38) (76.80,75.54) (78.90,83.06) (89.33,94.44) (71.72,62.21) (89.80,83.50) (60.73,73.03) (54.33,66.60) (64.08,71.64) (80.72,86.10) (74.39,73.35) (97.75,97.97) (79.18,91.95) (71.36,81.89) (85.41,84.19) (72.16,72.68) (77.27,56.87) (83.08,79.23) (69.10,60.24) (74.35,74.41) (75.37,74.80) (84.44,82.99) (83.80,68.06) (78.92,29.56) (69.34,70.87) (87.58,81.85) (83.19,69.66) (88.70,85.98) (80.89,91.37) (75.33,81.07) (57.48,74.24) (64.05,82.70) (73.33,68.70) (76.41,77.29) (80.38,78.99) (70.22,63.08) (63.84,71.86) (82.63,72.57) (58.44,70.68) (80.56,78.14) (82.54,87.40) (74.07,68.46) (84.27,87.12) (79.33,83.48) (73.29,71.53) (51.41,58.25) (94.31,89.50) (81.68,54.99)}
\newcommand{\ScatterTwo}{(84.52,80.86) (69.23,73.04) (70.61,70.28) (65.90,73.81) (82.44,78.29) (68.70,72.64) (77.37,78.49) (69.03,74.21) (71.07,67.85) (86.38,85.80) (75.99,87.61) (63.99,77.32) (62.50,61.15) (76.80,72.50) (78.90,79.78) (89.33,86.00) (71.72,69.51) (89.80,84.76) (60.73,63.59) (54.33,56.38) (64.08,63.48) (80.72,83.32) (74.39,66.46) (97.75,95.35) (79.18,91.55) (71.36,84.29) (85.41,86.10) (72.16,75.74) (77.27,76.41) (83.08,80.26) (69.10,68.72) (74.35,70.10) (75.37,67.00) (84.44,80.32) (83.80,73.44) (78.92,75.48) (69.34,70.22) (87.58,85.74) (83.19,80.10) (88.70,84.53) (80.89,88.18) (75.33,83.21) (57.48,73.05) (64.05,76.98) (73.33,77.44) (76.41,83.35) (80.38,77.17) (70.22,75.47) (63.84,56.34) (82.63,80.97) (58.44,65.23) (80.56,78.83) (82.54,76.70) (74.07,70.15) (84.27,86.29) (79.33,84.29) (73.29,72.07) (51.41,58.13) (94.31,77.94) (81.68,73.19)}
\newcommand{\ScatterSides}{60}
\newcommand{\BaseDecay}{0.0435}
\newcommand{\AppDecay}{0.14}
\newcommand{\CurveA}{103.1668}
\newcommand{\CurveC}{3.1669}
\newcommand{\WindowMin}{2}
\newcommand{\WindowMax}{8}
\newcommand{\WindowShare}{10}
\newcommand{\WeightMin}{0.5}
\newcommand{\WeightMax}{12}
\newcommand{\CurveBase}{(0,100.000) (0.5,98.778) (1,96.604) (1.5,94.477) (2,92.395) (2.5,90.359) (3,88.366) (3.5,86.416) (4,84.509) (4.5,82.642) (5,80.815) (5.5,79.028) (6,77.279) (6.5,75.568) (7,73.894) (7.5,72.256) (8,70.653) (8.5,69.085) (9,67.550) (9.5,66.049) (10,64.580) (10.5,63.142) (11,61.736) (11.5,60.360) (12,59.013) (12.5,57.695) (13,56.406) (13.5,55.145) (14,53.910) (14.5,52.703) (15,51.521) (15.5,50.365) (16,49.233) (16.5,48.126) (17,47.043) (17.5,45.983) (18,44.946) (18.5,43.932) (19,42.939) (19.5,41.967) (20,41.017) (20.5,40.087) (21,39.177) (21.5,38.286) (22,37.415) (22.5,36.563) (23,35.729) (23.5,34.912) (24,34.114) (24.5,33.332) (25,32.568) (25.5,31.820) (26,31.088) (26.5,30.372) (27,29.671) (27.5,28.985) (28,28.314) (28.5,27.658) (29,27.016) (29.5,26.387) (30,25.772) (30.5,25.170) (31,24.582) (31.5,24.005) (32,23.442) (32.5,22.890) (33,22.351) (33.5,21.823) (34,21.306) (34.5,20.800) (35,20.306) (35.5,19.822) (36,19.348) (36.5,18.885) (37,18.431) (37.5,17.988) (38,17.554) (38.5,17.129) (39,16.713) (39.5,16.307) (40,15.909)}
\newcommand{\CurveApp}{(0,100.000) (0.5,94.025) (1,87.522) (1.5,81.458) (2,75.805) (2.5,70.533) (3,65.618) (3.5,61.036) (4,56.763) (4.5,52.779) (5,49.064) (5.5,45.601) (6,42.371) (6.5,39.360) (7,36.553) (7.5,33.935) (8,31.494) (8.5,29.219) (9,27.097) (9.5,25.118) (10,23.274) (10.5,21.554) (11,19.950) (11.5,18.455) (12,17.061) (12.5,15.761) (13,14.549) (13.5,13.419) (14,12.365) (14.5,11.383) (15,10.467) (15.5,9.612) (16,8.816) (16.5,8.074) (17,7.381) (17.5,6.736) (18,6.134) (18.5,5.573) (19,5.049) (19.5,4.562) (20,4.107) (20.5,3.683) (21,3.287) (21.5,2.918) (22,2.575) (22.5,2.254) (23,1.955) (23.5,1.676) (24,1.417) (24.5,1.174) (25,0.948) (25.5,0.738) (26,0.541) (26.5,0.358) (27,0.188) (27.5,0.028) (28,0.000) (28.5,0.000) (29,0.000) (29.5,0.000) (30,0.000) (30.5,0.000) (31,0.000) (31.5,0.000) (32,0.000) (32.5,0.000) (33,0.000) (33.5,0.000) (34,0.000) (34.5,0.000) (35,0.000) (35.5,0.000) (36,0.000) (36.5,0.000) (37,0.000) (37.5,0.000) (38,0.000) (38.5,0.000) (39,0.000) (39.5,0.000) (40,0.000)}
\newcommand{\BaseAtTwo}{92}
\newcommand{\AppAtTwo}{76}
\newcommand{\BaseAtFive}{81}
\newcommand{\AppAtFive}{49}
\newcommand{\BaseAtTen}{65}
\newcommand{\AppAtTen}{23}
\newcommand{\DevMaeBefore}{8.0}
\newcommand{\DevMaeAfter}{4.3}
\newcommand{\DevCI}{[$-$4.90, $-$2.61]}
\newcommand{\DevRBefore}{0.61}
\newcommand{\DevRAfter}{0.86}
\newcommand{\DecayLow}{0.12}
\newcommand{\DecayHigh}{0.16}
\newcommand{\DecaySpread}{0.6}
\newcommand{\ExamplePlies}{27}
\newcommand{\ExampleSide}{Black}
\newcommand{\ExampleGraded}{11}
\newcommand{\ExampleWorst}{13\ldots{}h6}
\newcommand{\ExampleWorstLoss}{78.1}
\newcommand{\ExampleWorstBase}{1.3}
\newcommand{\ExampleWorstApp}{0.0}
\newcommand{\ExampleWeighted}{46.8}
\newcommand{\ExampleHarmonic}{12.3}
\newcommand{\ExampleBase}{29.6}
\newcommand{\ExampleApp}{75.5}
\newcommand{\ExampleChessCom}{78.9}
\newcommand{\ExampleAfter}{14.~\symqueen{}h7\#}
\newcommand{\ExampleGame}{% Written by projects/accuracy-vs-chesscom/paper/generate.mjs; do not edit by hand.
\begin{minipage}[t]{0.4\textwidth}\vspace{0pt}
\chessboard[setfen={r1bq1rk1/5ppp/p1n1p3/2bp2N1/5B2/3QP3/PPP2PPP/R4RK1 b - - 3 13}, boardfontsize=13pt, pgfstyle=straightmove, arrow=to, linewidth=0.07em, color=red!75!black, markmoves={h7-h6}, pgfstyle=straightmove, arrow=to, linewidth=0.07em, color=green!50!black, markmoves={f7-f5}]
\end{minipage}\hfill
\begin{minipage}[t]{0.56\textwidth}\vspace{0pt}\small
\textbf{The game:} 1.~d4 d5 2.~\symbishop{}f4 c5 3.~\symknight{}f3 e6 4.~e3 \symknight{}c6 5.~\symbishop{}b5 a6 6.~\symbishop{}a4 b5 7.~\symknight{}c3 bxa4 8.~\symknight{}xa4 \symknight{}f6 9.~dxc5 \symknight{}e4 10.~O-O \symknight{}xc5 11.~\symknight{}xc5 \symbishop{}xc5 12.~\symknight{}g5 O-O 13.~\symqueen{}d3 h6 14.~\symqueen{}h7\#\par
\textbf{Played:} 13\ldots{}h6 (red); the engine's best was 13\ldots{}f5 (green).\par
\textbf{UCI:} \uci{h7h6}; best \uci{f7f5}\par
\textbf{FEN:} \fen{r1bq1rk1/5ppp/p1n1p3/2bp2N1/5B2/3QP3/PPP2PPP/R4RK1 b - - 3 13}
\end{minipage}

\smallskip\small
\begin{tabular}{lrrr}
\toprule
Black's graded move & win\% lost & Lichess curve & steep curve \\
\midrule
3\ldots{}e6 & 0.4 & 99.4 & 95.8 \\
4\ldots{}\symknight{}c6 & 0.2 & 100.0 & 98.4 \\
5\ldots{}a6 & 2.1 & 91.9 & 74.6 \\
6\ldots{}b5 & 0.0 & 100.0 & 100.0 \\
7\ldots{}bxa4 & 0.3 & 99.5 & 96.3 \\
8\ldots{}\symknight{}f6 & 2.9 & 88.8 & 66.6 \\
9\ldots{}\symknight{}e4 & 6.4 & 75.9 & 39.9 \\
10\ldots{}\symknight{}xc5 & 0.0 & 100.0 & 100.0 \\
11\ldots{}\symbishop{}xc5 & 0.0 & 100.0 & 100.0 \\
12\ldots{}O-O & 3.8 & 85.4 & 58.7 \\
13\ldots{}h6 & 78.1 & 1.3 & 0.0 \\
\midrule
\multicolumn{2}{l}{the side's accuracy} & 29.6 & 75.5 \\
\bottomrule
\end{tabular}
}
\newcommand{\ReproTable}{%
\begin{tabular}{llrrrrrr}
\toprule
 & & \multicolumn{2}{c}{test MAE} & \multicolumn{2}{c}{train MAE} & \multicolumn{2}{c}{test $r$} \\
\cmidrule(lr){3-4}\cmidrule(lr){5-6}\cmidrule(lr){7-8}
evidence & instrument & recorded & rerun & recorded & rerun & recorded & rerun \\
\midrule
routine & baseline & 7.5 & 7.5 & 8.3 & 8.3 & 0.49 & 0.49 \\
routine & round 1 & 4.9 & 5.0 & 3.6 & 3.7 & 0.76 & 0.76 \\
routine & the app (v2) & 5.0 & 5.0 & 3.6 & 3.6 & 0.76 & 0.76 \\
phone & baseline & 7.5 & 7.6 & 8.1 & 8.2 & 0.50 & 0.51 \\
phone & round 1 & 5.4 & 5.5 & 4.6 & 4.5 & 0.72 & 0.71 \\
phone & the app (v2) & 5.4 & 5.5 & 4.6 & 4.5 & 0.72 & 0.71 \\
\bottomrule
\end{tabular}
}
\newcommand{\ReproLargest}{0.17}
\newcommand{\AppCommit}{\texttt{d8b7712}}
\newcommand{\RecordsCommit}{\texttt{b3b827d}}
\newcommand{\RereadRecordedMae}{0.34}
\newcommand{\RereadRecordedPNinety}{1.04}
\newcommand{\RereadRecordedSides}{454}
\newcommand{\RereadMae}{0.38}
\newcommand{\RereadPNinety}{1.16}
\newcommand{\RereadSides}{466}
\newcommand{\RereadMax}{7.0}
\newcommand{\RereadRuleMae}{0.5}
\newcommand{\RereadRuleP}{1.5}
\newcommand{\RereadHolds}{holds}
\newcommand{\RereadLaterRows}{42}
\newcommand{\RereadUntaggedSides}{286}
\newcommand{\RereadUntaggedMae}{0.03}
\newcommand{\RereadTaggedSides}{180}
\newcommand{\RereadTaggedMae}{0.92}


\title{Matching chess.com's accuracy\\[0.3em]
\large One round of hill-climbing a game review's accuracy, held out by game}
\author{Chess Coach research lab}
\date{}

\begin{document}
\maketitle

\begin{abstract}
The coaching app's game review gives each side an accuracy, and the player reads it beside the
one chess.com gives the same game. The app's first instrument followed Lichess's published
method: a score for each move from the share of winning chances it gives away, then a game
score that averages a volatility-weighted mean and a harmonic mean of those scores. On
\Games{} of the player's games, \GamesTest{} of them held out, it missed chess.com's number by
\BaseTestMae{} points on average, no better than guessing \ConstantGuess{} for every side
(\ConstantTestMae). One change, a plain mean of per-move scores on a steeper per-move curve,
cut the held-out error to \OneTestMae{} points (a paired change of \OneVsBase) and raised the
correlation with chess.com's number from \BaseTestR{} to \OneTestR. Merged with the app's
grading of moves an endgame tablebase decides, it is the instrument the app ships
(\TwoTestMae). The lab recomputes every number from the recorded per-game results, and reruns
all three instruments on engine evidence it rebuilt on the app as it now runs: no held-out
error moves by more than \ReproLargest{} points.
\end{abstract}

\section{The number the player compares}

After a game, chess.com shows each side an accuracy between 0 and 100, and so does the app's
game review. The player reads the two side by side, so a review whose accuracy disagrees with
chess.com's by a wide margin for no visible reason costs it the player's trust. chess.com's
formula is not published, and its number is not a ground truth about how well anyone played.
The goal here is narrower: parity with the number the player sees.

\section{Data and protocol}

\paragraph{Games.} The app's archive holds \Eligible{} of the player's games with chess.com's
accuracy for both sides, recorded when the daily routine fetched each game. \Games{} were
sampled, stratified by length (\ShortMax{} plies or fewer, \MediumMin{} to \MediumMax, more than \MediumMax) and
by result; the
other \DevPool{} form a dev pool. Each sampled game counts twice, once for each side.

\paragraph{Split.} The \Games{} games were split \GamesTrain{} for train and \GamesTest{} for
test, at random within each length bucket from a seeded generator, with the pilot games read
while building the eval kept in train. Seeds were tried in order until the baseline's
test-minus-train difference in both error and bias had a 95\% interval containing zero, so that
the split itself did not favour one half; \RejectedSplits{} seed was rejected. Work on a round saw
only the train games and the dev pool. The held-out games were scored once per round, for its
decision.

\paragraph{Evidence.} The app reviewed each game through its own entry point at two engine
budgets, and the documents were frozen: \emph{routine} (a sweep of \RoutineScan{} nodes a
position, every flagged move verified at \VerifyNodes{} nodes and \VerifyLines{} lines, as the
daily routine reviews a game), which gives the headline numbers, and \emph{phone} (a sweep of
\PhoneScan{} nodes, with verification for at most \PhoneShare\% of the moves, as the app reviews a
game on a phone), which is the guardrail. The engine is deterministic at a fixed node
count on one thread, so one run is the whole measurement, and an instrument is measured by
reading the frozen evidence again under its code.

\paragraph{Statistics.} The headline is the mean absolute error (MAE) between the app's
accuracy and chess.com's, over sides; within~\WithinPoints{} is the share of sides within \WithinPoints{}
points, the
bias is the mean signed error, and $r$ is Pearson's correlation over sides. Every interval is a
95\% bootstrap over games (a game's two sides share its evidence), and a variant's difference
from the baseline is paired on the same resampled games.

\section{The instruments}

Every instrument scores each graded move (book and forced moves are left out, and so are moves
without an exact evaluation, except in v2 those a tablebase decides) from $\Delta$, the share of the mover's winning chances the move
gives away against the engine's best, in percentage points:
\[
  s(\Delta) \;=\; \min\!\left(100,\ \max\!\left(0,\ \CurveA\, e^{-d\,\Delta} - \CurveC + 1\right)\right).
\]
\begin{itemize}[nosep]
\item \textbf{The baseline} is Lichess's published method: $d = \BaseDecay$, and a side's
accuracy is the average of two means of its scores, one weighted by volatility (the standard
deviation of the winning chances over a sliding window of positions, one position for every
\WindowShare{} of the side's graded moves and between \WindowMin{} and \WindowMax{} wide, each
weight held between \WeightMin{} and \WeightMax) and one harmonic.
\item \textbf{Round 1} keeps the curve's shape with a steeper decay, $d = \AppDecay$, and a
side's accuracy is the plain mean of its scores.
\item \textbf{The app (v2)} is round 1 over the moves the app now grades: since a later change,
a move an endgame tablebase decides is graded by its exact outcome.
\end{itemize}
A move that gives away 2 points of winning chances scores \BaseAtTwo{} on the baseline's curve
and \AppAtTwo{} on the steep one; 5 points, \BaseAtFive{} and \AppAtFive; 10 points,
\BaseAtTen{} and \AppAtTen{} (Figure~\ref{fig:curves}).

\begin{figure}[!htbp]
\centering
\begin{tikzpicture}
\begin{axis}[width=0.7\textwidth, height=5.2cm, xlabel={winning chances given away, $\Delta$ (points)},
  ylabel={move score}, xmin=0, xmax=40, ymin=0, ymax=102, grid=major, legend pos=north east,
  legend cell align=left]
\addplot[thick, black, dashed] coordinates {\CurveBase};
\addlegendentry{baseline, $d=\BaseDecay$}
\addplot[thick, red!70!black] coordinates {\CurveApp};
\addlegendentry{round 1 and the app, $d=\AppDecay$}
\end{axis}
\end{tikzpicture}
\caption{The per-move score on each curve, sampled from each instrument's own code.}
\label{fig:curves}
\end{figure}

\section{Why the baseline missed}

The baseline barely carried information about chess.com's number. Guessing the train sides'
mean chess.com accuracy, \ConstantGuess, for every side scores \ConstantTestMae{} on test and
\ConstantTrainMae{} on train, against the baseline's \BaseTestMae{} and \BaseTrainMae. The error
was not a constant offset (the baseline's bias on test is \BaseTestBias) but a tilt with the
game's length (Table~\ref{tab:tilt}): short games read about \TiltBaseShort{} points below
chess.com, long games about \TiltBaseLong{} above.

Both halves of the baseline's average explain the short end. In a short game a single blunder is
a large share of the side's moves; its score near zero sinks the harmonic mean, and its window
has the game's largest swing, so the volatility weights count it most. Figure~\ref{fig:example}
shows the held-out side the round improved most: \ExampleSide{} gives away little until
\ExampleWorst, which allows \ExampleAfter. Its score of \ExampleWorstBase{} pulls the harmonic mean
to \ExampleHarmonic{} and the weighted mean to \ExampleWeighted, and the side reads
\ExampleBase{} against chess.com's \ExampleChessCom. The plain mean of the same moves on the
steep curve reads \ExampleApp.

\begin{figure}[!htbp]
\ExampleGame
\caption{The held-out side whose error round 1 shrank most, from the evidence the lab rebuilt:
the worst of \ExampleSide's \ExampleGraded{} graded moves in a \ExamplePlies-ply game, the move
played (red) and the engine's best (green), with its score on both curves beside the side's
accuracy. chess.com gives the side \ExampleChessCom.}
\label{fig:example}
\end{figure}

\begin{table}[!htbp]
\centering\small
\TiltTable
\caption{Mean signed error (the app's accuracy less chess.com's) by the game's length, over all
\Games{} games' sides, routine evidence.}
\label{tab:tilt}
\end{table}

\section{Round 1: a plain mean on a steep curve}

\paragraph{Thesis.} chess.com's accuracy is close to an average of per-move scores, each
on a curve much steeper than Lichess's. The round's notes give the grounds: chess.com's
description of its method as mostly averaging, and the example points of an independent public
replication. A plain mean cannot be sunk by one move the way the harmonic mean is, and the
steeper curve makes up for it by charging small concessions more.

\paragraph{Before the held-out games.} On the dev pool (\DevPool{} games, never part of the
sample) the change took the MAE from \DevMaeBefore{} to \DevMaeAfter{} (paired 95\% interval
\DevCI) and $r$ from \DevRBefore{} to \DevRAfter. Decays from \DecayLow{} to \DecayHigh{} landed
within \DecaySpread{} points of each other, so the gain is the shape, not a tuned constant.

\paragraph{Result.} On the held-out games the MAE fell from \BaseTestMaeCI{} to \OneTestMaeCI,
a paired change of \OneVsBase{} points, and $r$ rose from \BaseTestRCI{} to \OneTestRCI{}
(Table~\ref{tab:results}, Figure~\ref{fig:scatter}). Train improved with it (\BaseTrainMae{} to
\OneTrainMae), and so did the phone guardrail (\PhoneBaseTestMae{} to \PhoneOneTestMae{} on
test). The length tilt shrinks but does not close: short games still read about
\TiltTwoShort{} points low and long ones \TiltTwoLong{} high.

\begin{table}[!htbp]
\centering\footnotesize
\setlength{\tabcolsep}{4pt}
\AggregateTable
\caption{Each instrument against chess.com, with 95\% bootstrap intervals over games. Every column
but the train MAE is on the held-out games. MAE and bias in accuracy points, the app's less
chess.com's; the last column is the change in test MAE from the baseline, paired on the same
games.}
\label{tab:results}
\end{table}

\begin{figure}[!htbp]
\centering
\begin{subfigure}{0.49\textwidth}\centering
\begin{tikzpicture}
\begin{axis}[scale only axis, width=0.78\linewidth, height=0.78\linewidth, xmin=20, xmax=100,
  ymin=20, ymax=100, xlabel={chess.com}, ylabel={the app}, grid=major]
\addplot[gray, densely dotted, domain=20:100] {x};
\addplot[only marks, mark=*, mark size=1.3pt, black] coordinates {\ScatterBase};
\end{axis}
\end{tikzpicture}
\caption{Baseline: MAE \BaseTestMae, $r = \BaseTestR$.}
\end{subfigure}\hfill
\begin{subfigure}{0.49\textwidth}\centering
\begin{tikzpicture}
\begin{axis}[scale only axis, width=0.78\linewidth, height=0.78\linewidth, xmin=20, xmax=100,
  ymin=20, ymax=100, xlabel={chess.com}, ylabel={the app}, grid=major]
\addplot[gray, densely dotted, domain=20:100] {x};
\addplot[only marks, mark=*, mark size=1.3pt, red!70!black] coordinates {\ScatterTwo};
\end{axis}
\end{tikzpicture}
\caption{The app (v2): MAE \TwoTestMae, $r = \TwoTestR$.}
\end{subfigure}
\caption{The \ScatterSides{} held-out sides, the app's accuracy against chess.com's, routine
evidence. The dotted line is parity.}
\label{fig:scatter}
\end{figure}

\paragraph{Merged with tablebase grading.} The app had meanwhile begun grading moves a tablebase
decides by their exact outcome. Round 1's formula over that graded set is the app's instrument,
v2, and costs nothing measurable: \TwoTestMae{} on test (paired change from the baseline
\TwoVsBase), \TwoTrainMae{} on train and \PhoneTwoTestMae{} on phone evidence.

\paragraph{What was tried and not run as rounds.} The hill-climb stopped after round 1 because
\StopRule. These were tried on the dev pool and train and not run as rounds:
\NotRun

\section{Moving the history}

The archive keeps every game's accuracy as the routine computed it, so a trend that mixed old
numbers with new ones would show a jump that is only the change of instrument. The app therefore
carries a stored accuracy only when it was taken under the current instrument, and otherwise reads
the game again from its stored per-move evaluations. That is safe only if those evaluations give
back the stored numbers under the old instrument, and the check of that was fixed before it ran:
a whole-game mean gap of at most \RereadRuleMae{} points and a 90th percentile of at most
\RereadRuleP. It held when recorded (\RereadRecordedMae{} and \RereadRecordedPNinety{} over
\RereadRecordedSides{} sides) and holds on the app's archive as it is now (\RereadMae{} and
\RereadPNinety{} over \RereadSides{} sides, the largest gap \RereadMax). Since the first run the
archive has gained \RereadLaterRows{} games stored under later instruments, which the old
instrument cannot reproduce and the check leaves out, as its definition says. The older rows,
stored with no instrument named, come back within \RereadUntaggedMae{} points on average
(\RereadUntaggedSides{} sides); the later rows named after the old instrument come back within
\RereadTaggedMae{} (\RereadTaggedSides{} sides). The rule does not separate the two, and why the
named rows reproduce less closely is not established here.

\section{Reproduction in the lab}

\paragraph{Recomputed.} The lab recomputes every aggregate from the per-game results the records
keep (\texttt{summarize.mjs}), and they \ReproAggregates. The records are read at \RecordsCommit.

\paragraph{Rerun.} The engine evidence the hill-climb used was lost, so the lab rebuilt it for
the same games on the app as it now runs (\AppCommit) and read it under each instrument
(\texttt{reproduce.mjs}). The app's review pipeline changed after the hill-climb, so the
evidence differs in detail, and the two older instruments needed one adjustment: they take the
root position's winning chances as an argument the current app no longer passes, so they run
with the app's document builder passing it again, as the app of their day did. Without it the
baseline's volatility weights fail on a game whose first moves are graded, and such a game reads
far below its recorded accuracy. With it, no held-out MAE moves by more than \ReproLargest{} points
(Table~\ref{tab:repro}).

\begin{table}[!htbp]
\centering\small
\ReproTable
\caption{The recorded numbers beside the lab's rerun on rebuilt evidence.}
\label{tab:repro}
\end{table}

\section{Limitations}

\begin{itemize}[nosep]
\item chess.com's formula is unknown, and parity with it is the goal, not a measure of play.
\item \Games{} games of one player, rated \PlayerRatingMin--\PlayerRatingMax{} on chess.com in
these games. The held-out half keeps the claim honest about this player, not about others.
\item The tilt with length remains: short games read low and long games high under every
variant tried.
\item The held-out games were scored for each round's decision. With one round kept and the
others stopped on the dev pool, the room for a lucky pick is small, but not zero.
\end{itemize}

\end{document}
